\section{Casos borde}

\subsection{Descripcion}

Este documento concentra los escenarios funcionales no triviales que pueden generar disputa, recalculo o intervencion administrativa.

\subsection{Alcance}

Aplica a partidos, resultados, lock, scoring y visibilidad.

\subsection{Definiciones}

\begin{itemize}
  \item \textbf{Suspendido}: partido interrumpido temporalmente sin resultado final utilizable.
  \item \textbf{Postergado}: partido movido a una nueva fecha.
  \item \textbf{Cancelado}: partido sin resultado oficial utilizable para la quiniela.
  \item \textbf{Corregido}: resultado oficial ajustado posteriormente.
\end{itemize}

\subsection{Reglas}

\begin{enumerate}
  \item Si un partido se posterga antes del lock, el lock se recalcula con el nuevo kickoff oficial.
  \item Si un partido se posterga despues del lock, los picks ya guardados permanecen validos y bloqueados, salvo decision administrativa extraordinaria auditada.
  \item Si un partido se suspende y luego se completa con resultado oficial compatible con la regla de 90 minutos mas agregado, el sistema puntua cuando ese resultado exista.
  \item Si un partido se cancela y no existe resultado oficial utilizable, no genera puntaje para ningun usuario.
  \item Si un resultado se corrige por fuente valida o override manual autorizado, se debe recalcular el puntaje afectado y dejar trazabilidad completa.
  \item Todo override manual de resultado debe registrar actor, momento, valor anterior, valor nuevo y justificacion.
\end{enumerate}

\subsection{Casos borde}

\begin{itemize}
  \item Si el proveedor de resultados entrega un valor temporal y luego otro definitivo, solo el resultado oficial vigente debe quedar activo.
  \item Si el sistema recalcula varias veces el mismo partido corregido, el proceso debe ser idempotente.
  \item Si un usuario intenta disputar un resultado, la evidencia operativa debe salir de la auditoria y de la fuente oficial aplicada.
\end{itemize}

\subsection{Invariantes}

\begin{itemize}
  \item Ninguna correccion manual debe borrar la historia del valor anterior.
  \item Ningun caso borde debe resolverse fuera de trazabilidad.
  \item Un partido cancelado no asigna puntos.
\end{itemize}

\subsection{Preguntas abiertas}

\begin{itemize}
  \item Confirmar la politica exacta de comunicacion al usuario cuando un resultado sea corregido despues de haberse mostrado.
\end{itemize}

\subsection{Decisiones pendientes}

\begin{itemize}
  \item `Decision pendiente` `Bloquea ambos`: definir el mensaje funcional al usuario para partidos corregidos.
  \item `Riesgo` `Bloquea backend`: confirmar criterios de prioridad entre API externa y override manual si ambos entregan cambios en ventanas cercanas.
\end{itemize}
